\begin{abstract}
Deployed language models nearly always carry persistent conditioning in the context
window, where it shares a channel with the interlocutor whose content it must resist.
We measure what that choice costs. A belief stated in context installs its stance in
98\% of generations but holds it under sustained adversarial pressure in only 58\%.
Compiling the same belief content into low-rank per-layer attention modulation (a
trained hypernetwork write function, in the manner of GenerativeAdapter and AdaLN-style
conditioning) inverts the trade-off: installation drops to 46\% while persistence rises
to 89\%. Composing the compiled channel with a small curated context slot recovers most
of the installation gap (77\%) and keeps the persistence advantage (89\%), a net gain of
31 percentage points in hold rate over context alone. The gap concentrates at logical
counterargument and emotional pressure, exactly where context-only conditioning
collapses to mean conviction 2.93 and 1.99 out of 5. Fixed-direction activation
steering extracted from identical belief content matches compiled modulation at
installation (42\% vs.\ 46\%) and collapses under pressure (14\% vs.\ 89\%), consistent
with an independent report that steering robustness degrades by up to 64 percentage
points under adversarial perturbation. The result replicates across three model
families. It is not a general anti-sycophancy trait: the update path concedes to
well-warranted evidence at \(2.4\times\) the rate of facially deficient evidence, with
0\% concession to contentless pressure, and general capability is unchanged within
margin (MMLU 76.5\% vs.\ 74.8\%, \(p=0.11\)). Installation and persistence are separable
properties of a conditioned belief, served by different channels, and a debate-oversight
protocol that places the instruction to hold in the same channel the opponent argues in
is targeting the wrong one.
\end{abstract}
